\chapter{Data flow and sequences}
\label{ch:flows}

This chapter follows the data through each major flow, in the order a user meets them. Each section
states what triggers the flow, what the client and the server do, and what the user sees. Sequence
diagrams use one lifeline per component; solid arrows are calls, dashed arrows are replies, and green
arrows are changes pushed by Realtime.

\section{Start-up and session restore}
\label{sec:flow-boot}

When the app starts, the root layout calls the store's \code{init()} once (\cref{fig:boot}). It registers
three listeners that last for the life of the process: an Auth listener that resets the app on sign-out,
an app-state listener that runs a catch-up when the app returns to the foreground, and a network listener
that runs a catch-up when the connection returns (\cref{sec:catch-up}).

\begin{figure}[ht]
  \centering
  % Start-up: restoring the session and choosing the first screen (Figure: start-up).
\begin{tikzpicture}[
    node distance=6mm and 9mm,
    st/.style={step=46mm, font=\sffamily\scriptsize},
    outcome/.style={terminal, font=\sffamily\scriptsize, text width=30mm, align=center},
    yes/.style={edgelabel, fill=none, font=\sffamily\scriptsize\itshape, text=muted},
  ]
  \node[st] (init) {Root layout calls \texttt{init()} once; it registers the auth, app-state and network listeners};
  \node[decision, below=of init] (session) {stored\\session?};
  \node[st, below=of session] (profile) {Load the profile row and sign its photos;\\load saved preferences};
  \node[decision, below=of profile] (loaded) {loaded?};
  \node[decision, below=of loaded] (exists) {profile\\exists?};
  \node[st, below=of exists] (extras) {Load the block list; restore the live session, chat or waiting request from \texttt{my\_status}};
  \node[st, below=of extras] (listen) {Subscribe to incoming requests; on the first join, look for one already waiting};

  \node[outcome, right=14mm of session] (welcome) {\texttt{(auth)} stack:\\welcome screen};
  \node[outcome, right=14mm of loaded] (cant) {\ui{Can't reach Yap}\\retry screen};
  \node[outcome, right=14mm of exists] (setup) {\texttt{profile-setup}\\stack};
  \node[outcome, right=14mm of listen] (app) {\texttt{(app)} stack:\\Go live tab};

  \draw[flow] (init) -- (session);
  \draw[flow] (session) -- node[yes, right] {yes} (profile);
  \draw[flow] (session) -- node[yes, above] {no} (welcome);
  \draw[flow] (profile) -- (loaded);
  \draw[flow] (loaded) -- node[yes, right] {yes} (exists);
  \draw[flow] (loaded) -- node[yes, above] {no (offline)} (cant);
  \draw[flow] (exists) -- node[yes, above] {no} (setup);
  \draw[flow] (exists) -- node[yes, right] {yes} (extras);
  \draw[flow] (extras) -- node[yes, right, align=left] {failures here are retried\\by the next catch-up} (listen);
  \draw[flow] (listen) -- (app);
  \draw[softflow] (cant.north) |- node[yes, pos=0.25, right, align=left] {retry, or a\\later catch-up} (profile.east);
\end{tikzpicture}

  \caption{Start-up. The store restores the session and the profile before any screen is shown, then the
    root guards choose the first stack.}
  \label{fig:boot}
\end{figure}

\code{init()} then restores the session that Supabase keeps in AsyncStorage. Without a session, the
guards show the welcome screen. With one, the store loads the user's profile row, signs its photo URLs
and loads the preferences saved on the phone. The profile is loaded before the session is published to
the screens, because a session without a profile means ``new account'' to the root layout, which would
otherwise flash profile setup. If the profile cannot be loaded, typically because the phone is offline at
launch, the store sets \code{loadFailed} and the root layout shows a retry screen instead.

For a user with a profile, the store loads the block list and calls \code{my_status} to restore the
state the app was in when it last closed: a live session (unless a request to go offline is still
waiting to reach the server), an open chat, or a request still waiting for its answer, which it starts
watching again. A failure in this step is tolerated and retried by the next catch-up. Finally the store
subscribes to incoming requests; when that channel first joins, it also looks for a request that arrived
before it was listening.

\section{Account creation, sign-in and password reset}
\label{sec:flow-account}

Yap uses Supabase Auth with email and password. Because the app does not open links, every email the
project sends carries an 8-digit code that the user types into the app (\cref{fig:seq-account}).

\begin{figure}[ht]
  \centering
  % Sequence: creating an account, confirming it and signing in (Figure: account sequence).
\begin{tikzpicture}
  \participant{ui}{0}{Auth screens}
  \participant[actor]{store}{3.5}{Store and\\API layer}
  \participant[serveractor]{auth}{7.0}{Supabase Auth}
  \participant[serveractor]{mail}{10.3}{SMTP sender\\(Resend)}
  \participant[serveractor]{db}{13.5}{Postgres}

  \fragment{ui}{db}{0.35}{7.35}{create an account}
  \msg{ui}{store}{1.0}{\code{signUp}(email, password)}
  \msg{store}{auth}{1.6}{\code{auth.signUp}}
  \msg{auth}{mail}{2.2}{confirmation email}
  \node[font=\sffamily\scriptsize, anchor=west, text=muted] at ($(mail |- 0,-2.2)+(0.15,0.12)$) {8-digit code};
  \msg[reply]{auth}{store}{2.8}{user, no session}
  \msg[reply]{store}{ui}{3.4}{open the code screen}
  \msg{ui}{store}{4.1}{\code{confirmSignUp}(email, code)}
  \msg{store}{auth}{4.7}{\code{verifyOtp}, type signup}
  \msg[reply]{auth}{store}{5.3}{session}
  \msg{store}{db}{5.9}{select own \code{profiles} row}
  \msg[reply]{db}{store}{6.5}{no row}
  \msg[reply]{store}{ui}{7.0}{guard opens profile setup}

  \fragment{ui}{db}{7.75}{12.55}{sign in}
  \msg{ui}{store}{8.4}{\code{signInWithPassword}}
  \msg{store}{auth}{9.0}{\code{auth.signInWithPassword}}
  \msg[reply]{auth}{store}{9.6}{session, or \code{email_not_confirmed}}
  \msg{store}{db}{10.2}{profile, block list, \code{my_status}}
  \msg[reply]{db}{store}{10.8}{rows}
  \msg[reply]{store}{ui}{11.4}{guard opens the app}
  \seqnote{ui}{store}{11.75}{unconfirmed: send a new code, open the code screen}

  \lifelines{12.9}{ui, store, auth, mail, db}
\end{tikzpicture}

  \caption{Creating an account, confirming it with the emailed code, and signing in.}
  \label{fig:seq-account}
\end{figure}

\paragraph{Creating an account} The sign-up screen checks that the password has at least eight
characters and matches its confirmation, then calls \code{signUp}. Supabase creates the user and emails a
confirmation code. If the email already belongs to a confirmed account, Supabase returns a user with no
identities and sends nothing. The app opens the code screen exactly as for a new email, and
\ui{Resend code} behaves the same too, so sign-up never shows whether someone has an account; a note
under the code says \ui{No code? If you already have an account, sign in or reset your password}, with
both as links. If \ui{Confirm email} were
switched off in the dashboard, sign-up would return a session at once and the code screen would be
skipped.

\paragraph{Confirming the email} The code screen accepts exactly eight digits through one hidden input
behind eight boxes, which keeps the operating system's one-time-code autofill working. Because the input
is hidden, \ui{Paste code} (or a long press on the boxes) reads the clipboard, keeps the eight-digit code
from whatever was copied, and fills it in; if there is no code, it says so. The emails show the code on
its own large line so it is easy to copy. The screen submits as soon as the eighth digit arrives. Verifying the code with type \code{signup} returns a
session; the store then loads the (missing) profile and the guards open profile setup. \ui{Resend code}
is disabled for 60 seconds after each send, matching the hosted project's limit of one code per address
per minute (\cref{tab:dashboard}). The local stack allows one per second, so the button's wait is the
only limit there.

\paragraph{Signing in} Sign-in calls \code{signInWithPassword}. If the password is correct but the
account was never confirmed, Supabase answers \code{email_not_confirmed}; the screen then requests a new
code and opens the code screen. Errors are reported as sentences: a wrong password, for example, suggests
\ui{Forgot password}, because accounts created before passwords existed have none.

\paragraph{Resetting a password} \Cref{fig:seq-reset} shows the reset flow. \ui{Forgot password} emails
an 8-digit reset code. The code screen, opened with purpose \code{reset}, only collects the code: it is
checked when the new password is saved, because verifying a recovery code signs the user in. On
\ui{Save and sign in}, the store verifies the code with type \code{recovery}, which creates a session, and
saves the new password with \code{updateUser}. If saving fails, the store signs out locally, so the
recovery session cannot sign the user in on the next launch without a password, and says that the code
has been used up. A signed-in user can also change the password from Settings, which calls
\code{updateUser} directly.

\begin{figure}[ht]
  \centering
  % Sequence: resetting a forgotten password with an emailed code (Figure: reset sequence).
\begin{tikzpicture}
  \participant{ui}{0}{Forgot password,\\code, new password}
  \participant[actor]{store}{4.2}{Store and\\API layer}
  \participant[serveractor]{auth}{8.6}{Supabase Auth}
  \participant[serveractor]{mail}{12.8}{SMTP sender\\(Resend)}

  \msg{ui}{store}{0.8}{\code{requestPasswordReset}(email)}
  \msg{store}{auth}{1.4}{\code{resetPasswordForEmail}}
  \msg{auth}{mail}{2.0}{reset email with an 8-digit code}
  \seqnote{ui}{ui}{2.35}{the code screen keeps the code;\\it is not checked yet}
  \msg{ui}{store}{3.6}{\code{finishPasswordReset}(email, code, password)}
  \msg{store}{auth}{4.2}{\code{verifyOtp}, type recovery}
  \msg[reply]{auth}{store}{4.8}{session (signed in)}
  \msg{store}{auth}{5.4}{\code{updateUser}(password)}
  \msg[reply]{auth}{store}{6.0}{saved}
  \msg[reply]{store}{ui}{6.6}{load profile; guard opens the app}
  \fragment{ui}{auth}{7.0}{8.5}{if saving fails}
  \selfmsg{store}{7.55}{\code{signOut}, local scope}
  \msg[reply]{store}{ui}{8.2}{``the code is used up''}

  \lifelines{8.75}{ui, store, auth, mail}
\end{tikzpicture}

  \caption{Resetting a forgotten password. The emailed code is checked only together with the new
    password.}
  \label{fig:seq-reset}
\end{figure}

\section{Profile setup and saving}
\label{sec:flow-profile}

A signed-in user without a profile row goes through seven screens: a greeting with the 18+ confirmation,
then six numbered steps for the name, gender, introduction, interests, profile picture and further
photos. The answers are kept in the profile draft until the last step, \ui{Finish}, calls the store's
\code{saveProfile}. Editing the profile later uses the same action from a single-page form.
\Cref{fig:seq-profile} shows what saving does.

\begin{figure}[ht]
  \centering
  % Sequence: saving a profile and uploading its photos (Figure: profile sequence).
\begin{tikzpicture}
  \participant{ui}{0}{Setup or edit\\screen}
  \participant[actor]{store}{3.0}{Store}
  \participant[actor]{api}{6.2}{API layer}
  \participant[serveractor]{storage}{9.6}{Storage\\(photos bucket)}
  \participant[serveractor]{db}{13.0}{Postgres}

  \msg{ui}{store}{0.8}{\code{saveProfile}(values)}
  \msg{store}{api}{1.4}{\code{api.saveProfile}}
  \fragment{api}{storage}{1.75}{4.4}{for each new photo}
  \selfmsg{api}{2.35}{re-encode as JPEG, at most\\1080\,px wide, quality 0.8}
  \msg{api}{storage}{3.45}{upload \code{<id>/<slot>-<ms>.jpg}}
  \msg{api}{storage}{4.05}{profile picture: also a 320\,px copy}
  \msg{api}{db}{4.90}{does my row exist?}
  \msg{api}{db}{5.50}{insert (first save) or update editable columns}
  \seqnote{db}{db}{5.80}{check constraints: own folder,\\at most 6 extra photos}
  \msg{api}{storage}{6.90}{list own folder; remove files no longer used}
  \msg{api}{db}{7.50}{select own row}
  \msg{api}{storage}{8.10}{sign photo URLs}
  \msg[reply]{api}{store}{8.70}{stored profile}
  \selfmsg{store}{9.10}{set \code{profile}; on the first save, mark the\\tutorial due and listen for requests}
  \msg[reply]{store}{ui}{10.10}{guard opens the app}

  \lifelines{10.45}{ui, store, api, storage, db}
\end{tikzpicture}

  \caption{Saving a profile. Photos are uploaded before the row is written, so the row only ever points
    at files that exist.}
  \label{fig:seq-profile}
\end{figure}

For each photo, the API layer decides whether it is new. A \code{file://} address comes from the picker
or the camera: it is re-encoded as a JPEG at most 1080 pixels wide (never enlarged) at quality 0.8 and
uploaded to \code{<user-id>/<slot>-<milliseconds>.jpg}. The profile picture is also uploaded as a
320-pixel copy for avatars; if that upload fails, avatars fall back to the full photo. An
\code{https://} address is a signed URL of a photo already stored, so its storage path is recovered from
the URL and kept.

The API layer then inserts the row on the first save, or updates only the editable columns afterwards.
It lists the user's folder and removes any file the saved profile no longer uses, such as replaced photos
or uploads from a failed attempt; this tidy-up is best-effort and is retried by the next successful save.
Finally it reloads the row and signs its photo URLs. On the first save, the store marks the tutorial as
due, starts listening for incoming requests, and the guards open the app.

\section{Going live and the nearby list}
\label{sec:flow-live}

\Cref{fig:seq-live} shows the life of a live session from the client's point of view.

\begin{figure}[ht]
  \centering
  % Sequence: going live, refreshing the nearby list and going offline (Figure: live sequence).
\begin{tikzpicture}
  \participant{ui}{0}{Screens and\\app layout}
  \participant[actor]{loc}{3.0}{Phone location}
  \participant[actor]{store}{6.2}{Store and\\API layer}
  \participant[serveractor]{db}{9.8}{Postgres}
  \participant[serveractor]{storage}{13.3}{Storage}

  \msg{ui}{loc}{0.8}{ask for foreground permission}
  \msg[reply]{loc}{ui}{1.4}{position and accuracy}
  \seqnote{loc}{loc}{1.65}{fresh fix within 8\,s, else\\one at most 5 min old}
  \msg{ui}{store}{2.9}{\code{goGreen}(minutes, note, position)}
  \msg{store}{db}{3.5}{\code{go_green}}
  \seqnote{db}{db}{3.8}{relocation rules;\\insert or update the session}
  \msg[reply]{db}{store}{4.9}{server start and end times}
  \msg[reply]{store}{ui}{5.5}{live: button turns red, screen kept awake, end timer set}

  \fragment{ui}{storage}{5.95}{9.4}{loop: every 20 s while live and in the foreground}
  \msg{ui}{store}{6.8}{\code{refreshNearby}}
  \msg{store}{db}{7.4}{\code{nearby}}
  \msg[reply]{db}{store}{8.0}{up to 50 person cards}
  \msg{store}{storage}{8.6}{sign profile pictures (cached)}
  \msg[reply]{store}{ui}{9.15}{nearby list; prompt for new people}

  \fragment{ui}{storage}{9.8}{12.8}{go offline (tap the red button), or the end time is reached}
  \msg{ui}{store}{10.65}{\code{goOffline}}
  \selfmsg{store}{10.9}{clear the live state at once}
  \msg{store}{db}{11.85}{\code{go_offline}: withdraw waiting request, delete session}
  \seqnote{store}{db}{12.15}{on failure: marked pending, retried by the next catch-up}

  \lifelines{13.05}{ui, loc, store, db, storage}
\end{tikzpicture}

  \caption{Going live, keeping the nearby list fresh, and going offline.}
  \label{fig:seq-live}
\end{figure}

\paragraph{Going live} Tapping the green \ui{YAP} button opens the go-live sheet, where the user chooses
30, 60 or 120 minutes and may write a note. On \ui{Go live}, the app asks for foreground location
permission. It never asks for background location. It requests a position with balanced accuracy and
waits up to eight seconds; if no fix arrives, it uses a position the phone recorded in the last five
minutes. The store sends the position, its accuracy, the duration and the note to \code{go_green}, which
applies the relocation rules and writes the session (\cref{sec:relocation-limits}). The function returns
the session's start and end times as the server set them, so the countdown on the button ends when the
session really does, whatever the phone's clock says.

While the user is live, the \code{(app)} layout keeps the screen awake, sets a timer that takes the user
offline at the end time, and polls the nearby list. The button turns red, reads \ui{LIVE} with the
minutes left, and a tap takes the user offline.

\paragraph{The nearby list} One shared list in the store serves the \ui{Nearby} tab, the person sheet and
the nearby prompts. A single poll refreshes it every 20 seconds while the user is live and the app is in
the foreground, and pulling the list down refreshes it at once. Each refresh calls \code{nearby}, which
returns person cards and nothing about where anyone is. The API layer signs the profile pictures and their
small copies in one batch; a person's further photos are signed only when their profile is opened. A
counter discards an answer that arrives after a newer refresh has started.

\paragraph{Nearby prompts} The hook \code{useNearbyAlerts} compares each new list with the people already
announced in this live session. On the first list, if more than one person is there, it shows a summary
(``3 people near you are up for a chat''); after that it announces a newly arrived person by name, the
most recent one if several arrive in the same refresh. Each prompt plays a haptic and disappears after
eight seconds, and each person is announced at most once per session.

\paragraph{Going offline} A tap on the red button, or the end-time timer, calls \code{goOffline}. The store
clears the live state at once and then calls \code{go_offline}, which withdraws the user's waiting request
and deletes the session row. If the call fails when the user chose to go offline, the store records that
the request is pending, shows \ui{Almost offline}, and retries at the next catch-up. If a retry is in
progress when the user goes live again, \code{goGreen} waits for it to finish, so the retry cannot delete
the new session. A session that simply runs out is
also removed by the server's minute job, so a failed call at the end time needs no retry.

\section{Saying hi, answering and chatting}
\label{sec:flow-sayhi}

\subsection{Sending a request and receiving the answer}

A user opens a person from the nearby list or a prompt and taps \ui{Say hi}. The store refuses locally
while the user is in a chat or already has a request waiting; the server applies the same rules and more
(\cref{tab:functions}). \Cref{fig:seq-request} shows the exchange that follows.

\begin{figure}[ht]
  \centering
  % Sequence: a say-hi request and its answer (Figure: request sequence).
\begin{tikzpicture}
  \participant{a}{0}{Sender's app}
  \participant[serveractor]{db}{3.9}{Postgres}
  \participant[serveractor]{rt}{7.8}{Realtime}
  \participant{b}{11.6}{Recipient's app}

  \msg{a}{db}{0.8}{\code{send_request}(target)}
  \seqnote{db}{db}{1.05}{checks; insert \code{pending},\\ends in 5 min}
  \msg[reply]{db}{a}{2.1}{request ID and end time}
  \selfmsg{a}{2.35}{wait: listen for updates,\\poll every 7\,s, give up 3\,s after the end}
  \msg[push]{db}{rt}{3.4}{row inserted}
  \msg[push]{rt}{b}{4.0}{insert where \code{to_user} = recipient}
  \msg{b}{db}{4.6}{\code{request_person}}
  \msg[reply]{db}{b}{5.2}{sender's person card}
  \selfmsg{b}{5.45}{full-screen request,\\5-minute countdown}

  \fragment{a}{b}{6.4}{14.0}{alt}
  \fragguardat{a}{6.75}{[recipient accepts]}
  \msg{b}{db}{7.35}{\code{respond}(accept)}
  \seqnote{db}{db}{7.6}{\code{accepted}; both users \code{chatting}}
  \msg[push]{db}{rt}{8.4}{row updated}
  \msg[push]{rt}{a}{9.0}{status \code{accepted}}
  \seqnote{a}{a}{9.25}{both apps open the chat}
  \fragsep{a}{b}{10.0}{[recipient declines, or nobody answers in time]}
  \msg{b}{db}{10.65}{\code{respond}(decline), or \code{expire_stale}}
  \msg[push]{db}{rt}{11.25}{row updated}
  \msg[push]{rt}{a}{11.85}{status \code{declined} or \code{expired}}
  \seqnote{a}{a}{12.1}{``Not this time''}
  \fragsep{a}{b}{12.85}{[sender cancels]}
  \msg{a}{db}{13.7}{\code{cancel_request}: cancels it, even if just accepted}

  \lifelines{14.25}{a, db, rt, b}
\end{tikzpicture}

  \caption{A say-hi request. The sender learns the outcome from Realtime, from a slow poll, or from its
    own timer, whichever comes first. The server refuses an acceptance after the end time; the timer's
    three seconds of grace only wait for news of an earlier acceptance that is still on its way.}
  \label{fig:seq-request}
\end{figure}

\code{send_request} returns the request's ID and end time, and the sender's app opens the waiting
screen with a five-minute countdown. The store then waits for the answer in three ways at once, so that
one of them always works: a Realtime channel for updates to that request row, a query of its status every
seven seconds, and a timer that gives up three seconds after the end time. The first of these to report
an outcome settles the wait: \code{accepted} opens the chat; any other final status, or the timer, shows
\ui{Not this time}, which reads the same whether the recipient declined or never answered. Separately,
the waiting screen's own countdown shows \ui{Not this time} as soon as the end time arrives. The recipient
cannot accept after the end time: \code{respond} then refuses with \code{request expired}. The three
seconds of grace only wait for news of an acceptance made just before the end time, whose Realtime update
or poll result is still on its way; if it arrives in that window, the chat replaces \ui{Not this time}.
If the sender has already tapped \ui{Back to nearby list}, the store has dropped the request, so the
wait ignores the answer and the next catch-up opens the chat instead.

On the recipient's phone, the insert reaches the app through the incoming-requests channel, which
filters on \code{to_user}. The store fetches the sender's card with \code{request_person} and the
\code{(app)} layout opens a full-screen request with the sender's profile, a countdown, and two buttons of
equal size: \ui{Not now} and \ui{Accept}. While a request is already on screen or the user is in a chat,
later requests are not shown; they expire on the server, and their senders see \ui{Not this time}. If the
recipient's app was in the background or offline when the request arrived, the catch-up finds it by
querying for the newest waiting request addressed to the user.

\paragraph{Accepting} \code{respond} with \code{accept = true} sets the request to \code{accepted} and both
users to \code{chatting}, which removes them from everyone else's nearby list. If the recipient had a
request of their own waiting, the server withdraws it and the app stops showing it. The recipient's app
opens the chat; the sender's app opens it when the update arrives. If the accept fails, for example
because the sender cancelled a moment earlier or is no longer live, the recipient sees a short
explanation.

\paragraph{Cancelling} \ui{Cancel request} on the waiting screen removes the request from the screen at
once and calls \code{cancel_request}. If that call fails, the request reappears, the app resumes watching
for the answer, and an alert explains that the user is still waiting. If the cancel reaches the server
just after an accept, the server ends that chat too, so nobody is left in a chat the other person never
sees.

\subsection{The one-to-one chat}

Both apps open the chat in the same way (\cref{fig:seq-chat}): the store subscribes to the chat's
channel first and loads the newest 200 messages each time the channel joins, so a message sent between
the two steps is not lost. The chat screen shows the partner's photo and name, a safety tip during the user's first few chats, and
a message composer. Tapping the partner's name opens their profile, including the live note that often
says where to find them.

\begin{figure}[ht]
  \centering
  % Sequence: the one-to-one chat, from opening to ending (Figure: chat sequence).
\begin{tikzpicture}
  \participant{a}{0}{User A's app}
  \participant[serveractor]{db}{4.3}{Postgres}
  \participant[serveractor]{rt}{8.6}{Realtime}
  \participant{b}{12.9}{User B's app}

  \msg{a}{rt}{0.8}{join \code{chat-<id>}: messages, request updates}
  \msg[reply]{rt}{a}{1.4}{subscribed}
  \msg{a}{db}{2.0}{newest 200 messages; request status}
  \msg[reply]{db}{a}{2.6}{history}
  \seqnote{a}{b}{2.95}{B's app opens the chat in the same way}
  \msg{a}{db}{4.0}{\code{send_chat_message}(id, text)}
  \msg[reply]{db}{a}{4.6}{saved row, shown at once}
  \msg[push]{db}{rt}{5.2}{row inserted}
  \msg[push]{rt}{b}{5.8}{new message}
  \msg[push]{rt}{a}{6.4}{same message, skipped by ID}
  \msg{b}{db}{7.2}{\code{end_chat}(id, met or cancelled)}
  \seqnote{db}{db}{7.45}{status set; both users\\\code{open} again}
  \msg[push]{db}{rt}{8.5}{row updated}
  \msg[push]{rt}{a}{9.1}{status no longer \code{accepted}}
  \selfmsg{a}{9.35}{``B ended the chat'';\\composer closed}
  \seqnote{a}{b}{10.35}{a block ends the chat in the same way; if B deletes the account, A's app notices at its next catch-up or send}

  \lifelines{10.85}{a, db, rt, b}
\end{tikzpicture}

  \caption{The one-to-one chat. Messages are merged by ID, so the copy that Realtime delivers to the
    sender is ignored.}
  \label{fig:seq-chat}
\end{figure}

\code{send_chat_message} returns the saved row, which the sender's app shows immediately; when Realtime
delivers the same row a moment later it is recognised by its ID and skipped. Messages are kept sorted by
the time the server stored them.

\paragraph{Ending the chat} \ui{End} asks \ui{Did you meet?} with three choices: \ui{Keep chatting},
\ui{We met} and \ui{End chat}. The two ending choices close the chat on this phone at once and call
\code{end_chat} with outcome \code{met} or \code{cancelled}; both users return to the nearby lists if they
are still live. The other phone learns of the end through the chat channel, which also listens for
updates to the request row: when the status leaves \code{accepted}, the chat stays on screen with the
notice that the partner ended it and the composer replaced by that notice. Leaving the chat screen without
ending it keeps the chat open; a bar on the \ui{Go live} and \ui{Nearby} tabs reopens it.

\paragraph{Chats nobody ends} If both users simply stop using the app, the minute job ends the chat once
neither of them is live (\cref{sec:housekeeping}), and \code{my_status} no longer offers it after a
restart.

\section{Events and group chats}
\label{sec:flow-events}

\Cref{fig:seq-events} traces the event features.

\begin{figure}[ht]
  \centering
  % Sequence: loading, creating and joining events, and the group chat (Figure: events sequence).
\begin{tikzpicture}
  \participant{ui}{0}{Map and\\event screens}
  \participant[actor]{loc}{3.2}{Phone location}
  \participant[actor]{store}{6.5}{Store and\\API layer}
  \participant[serveractor]{db}{10.1}{Postgres}
  \participant[serveractor]{rt}{13.6}{Realtime}

  \msg{ui}{loc}{0.8}{permission (asked once) and position}
  \msg[reply]{loc}{ui}{1.4}{recent, fresh or last known position}
  \msg{ui}{store}{2.0}{\code{loadEvents}(centre)}
  \msg{store}{db}{2.6}{\code{events_near}(lat, lng, radius 1--50\,km)}
  \msg[reply]{db}{store}{3.2}{up to 50 event cards}
  \msg[reply]{store}{ui}{3.8}{pins and the radius circle}

  \fragment{ui}{rt}{4.2}{6.8}{create: long-press the map, or ``Create event here''}
  \msg{ui}{store}{5.05}{\code{createEvent}(category, title, details, place, start)}
  \msg{store}{db}{5.65}{\code{create_event}: host becomes a member}
  \msg[reply]{db}{store}{6.25}{event card; the group chat opens}

  \fragment{ui}{rt}{7.2}{10.4}{join: tap a pin, then ``Join and open chat''}
  \msg{ui}{loc}{8.05}{position}
  \msg{ui}{store}{8.65}{\code{joinEvent}(id, position)}
  \msg{store}{db}{9.25}{\code{join_event}: within 5\,km}
  \msg[reply]{db}{store}{9.85}{member; the group chat opens}

  \fragment{ui}{rt}{10.8}{13.9}{group chat, while its screen is open}
  \msg{store}{rt}{11.65}{join \code{messages-event_id-<id>}}
  \msg{store}{db}{12.25}{newest 200 messages; \code{send_event_message}}
  \msg[push]{rt}{store}{12.85}{new messages from members}
  \msg[reply]{store}{ui}{13.45}{thread, minus blocked senders}

  \lifelines{14.15}{ui, loc, store, db, rt}
\end{tikzpicture}

  \caption{Loading events around the user, creating, joining and deleting an event, and the group chat.}
  \label{fig:seq-events}
\end{figure}

\paragraph{Finding events} When the \ui{Map} tab opens, and each time its tab button is pressed, the
screen finds the user. It asks for location permission only if the user has never decided, uses a
position the phone recorded in the last two minutes if there is one, otherwise waits up to eight seconds
for a fresh fix. If neither comes, it only centres the map on the last position the phone knows, however
old, and shows that it could not find the user: an old position is never sent to the server, which would
otherwise hold it as the user's last known position and refuse the real one. With a recent or fresh
position it centres the map and loads the events within the chosen radius of that point: the store sends
the position to \code{set_map_position} and then calls \code{events_near}, which lists events around the
position the server holds (\cref{sec:location-privacy}). If the server refuses the new position because it is too
far from the last known one or changed too often, the map keeps the pins it has and shows the refusal
under the banner with a \ui{Try again} button. Without a position, the
map loads no events and shows a note and a button to turn location on; when the user returns from the
phone's settings, the map tries again.

The radius is chosen with a slider from 1 to 50 kilometres (default 10) and drawn on the map as a circle.
Releasing the slider saves the value for the account, fits the map to the circle, and lists the events
again with \code{events_near} alone, around the position the server already holds: no position is sent
again, so one the phone found a while ago can never replace a newer one. \code{events_near} applies the radius on the server and never searches beyond 50
kilometres. Each event becomes a pin showing its category icon and title; events the user has joined are
drawn in green. If none are in range, a note suggests widening the distance or creating an event.

\paragraph{Creating an event} A long press on the map, or \ui{Create event here} for the centre of the
current view, opens the event form at that place. The host picks one of twelve categories, writes a title
and optional details, and picks a start time from now to seven days ahead; on Android the picker asks for
the date and then the time. \code{create_event} stores the event, adds the host as a member and returns
the event card, and the app opens the event's group chat.

\paragraph{Joining} Tapping a pin opens the event sheet with the category, title, start time
(\ui{Starts in 25 min}, \ui{Started 10 min ago}), number of members, details and host. \ui{Join and open
chat} requests the phone's position, sends it to \code{set_map_position}, and calls \code{join_event},
which accepts the user only within five kilometres of the position the server holds. A user therefore
sees events up to 50 kilometres from their own position, but only someone nearby can join one and read
its chat. The host's photo is sent only to people who joined, so after a join the app lists the pins
again with \code{events_near} alone to show it. If the host has deleted the event in the meantime, the server
answers \code{no such event}; the app takes the event off the map and the sheet shows \ui{This event was
deleted}.

\paragraph{The group chat} The group chat subscribes to the event's messages while its screen is open and
loads the newest 200 messages each time the channel joins. Messages show the sender's name. Messages from
users the reader has blocked are excluded by row-level security and filtered again on screen. A long press
on someone else's message offers \ui{Report} and \ui{Block}; blocking the host also removes the event from
the map and closes the chat. A guest can leave the event, which removes their membership and their access
to its chat. The host has no \ui{Leave} button, because on the server a host leaving deletes the event;
the host deletes it deliberately instead. Sending to the chat of an event that no longer exists, or that
the user has left, is refused with \code{not a member}, shown as \ui{You can't post here. The event may
have ended or been deleted.}

\paragraph{Deleting an event} The event sheet shows the host a \ui{Delete event} button. After a
confirmation (\ui{Delete this event?} \ui{Everyone who joined will lose the group chat.}), the store
calls \code{delete_event}, which accepts only the host. Deleting the row removes the event's members and
its group-chat messages through their cascading foreign keys. The app takes the event off the map at once;
other phones stop showing it at their next load of events. An event already deleted, for example from
another phone, counts as deleted.

\section{Photos}
\label{sec:flow-photos}

Uploading photos is part of saving a profile (\cref{sec:flow-profile}). Showing another user's photo
takes a signed URL, because the bucket is private (\cref{fig:seq-photos}).

\begin{figure}[ht]
  \centering
  % Sequence: showing another user's photo through a signed URL (Figure: photo sequence).
\begin{tikzpicture}
  \participant{ui}{0}{List or\\profile screen}
  \participant[actor]{api}{3.3}{API layer\\(URL cache)}
  \participant[serveractor]{storage}{6.9}{Storage}
  \participant[serveractor]{db}{10.4}{Postgres\\read policy}
  \participant[actor]{img}{13.8}{Image cache}

  \msg{ui}{api}{0.8}{photo paths of the cards}
  \selfmsg{api}{1.05}{skip paths with a URL that has 10 min\\left, or a refusal under 1 min old}
  \msg{api}{storage}{2.1}{\code{createSignedUrls}(the rest, 1 h)}
  \msg{storage}{db}{2.7}{select policy \code{photos_read_visible}}
  \selfmsg{db}{2.95}{\code{can_see_photo}\\for each file}
  \msg[reply]{db}{storage}{3.9}{allowed or not}
  \msg[reply]{storage}{api}{4.5}{a URL for each allowed path}
  \selfmsg{api}{4.75}{cache: a URL for 1 h,\\a refusal for 60\,s}
  \msg[reply]{api}{ui}{5.7}{display URLs}
  \msg{ui}{img}{6.3}{show, with the token removed from the cache key}
  \msg{img}{storage}{6.9}{download only on a cache miss}

  \lifelines{7.25}{ui, api, storage, db, img}
\end{tikzpicture}

  \caption{Showing another user's photo. Storage signs a URL only if the read policy allows the signed-in
    user to see the file.}
  \label{fig:seq-photos}
\end{figure}

The API layer keeps a cache from storage path to signed URL. A URL is created for fifteen minutes and
reused until fewer than five minutes remain, so the same photo keeps the same URL across refreshes. A path that
could not be signed, because the file does not exist or its owner is not visible to the user, is
remembered for one minute and then asked for again, so the photo appears soon after its owner becomes
visible. Expired entries are dropped whenever new URLs are signed, and the whole cache is cleared on
sign-out. All the paths a screen needs are signed in one request.

Storage signs a URL only if the user may read the object. The read policy calls
\code{private.can_see_photo}, which allows a photo only while its owner is visible to the user
(\cref{sec:photo-privacy}). On the device, \code{expo-image} caches the downloaded image under the URL
without its token, so re-signing a URL does not cause the image to be downloaded again.

\section{Account deletion}
\label{sec:flow-delete}

\ui{Delete account} in Settings opens a screen that lists what is deleted. Its \ui{Delete my account}
button only works once the user has typed their first name as it is on their profile (ignoring capitals
and spaces around it), so a stray tap cannot delete an account. The app checks for a connection first,
so a dropped network cannot remove the photos and leave the account, and a second tap waits for the
first. It then runs the sequence in \cref{fig:seq-delete}. The photos are removed first, through the Storage API, because once the account is
gone the session can no longer reach Storage, and Supabase does not allow deleting rows from the storage
tables directly.

\begin{figure}[ht]
  \centering
  % Sequence: deleting an account (Figure: deletion sequence).
\begin{tikzpicture}
  \participant{ui}{0}{Settings}
  \participant[actor]{store}{3.5}{Store and\\API layer}
  \participant[serveractor]{storage}{6.9}{Storage}
  \participant[serveractor]{db}{10.2}{Postgres}
  \participant[serveractor]{auth}{13.6}{Supabase Auth}

  \selfmsg{ui}{0.6}{confirm:\\``Delete your account?''}
  \msg{ui}{store}{1.6}{\code{deleteAccount}}
  \msg{store}{storage}{2.2}{list own folder, 100 at a time}
  \msg[reply]{storage}{store}{2.8}{paths}
  \msg{store}{storage}{3.4}{remove every file}
  \msg{store}{db}{4.0}{\code{delete_account}}
  \selfmsg{db}{4.25}{chat partner back to \code{open};\\delete the \code{auth.users} row;\\every user table cascades}
  \msg[reply]{db}{store}{5.6}{done}
  \msg{store}{auth}{6.2}{\code{signOut}}
  \selfmsg{store}{6.45}{auth listener: close channels,\\forget URLs, reset state}
  \msg[reply]{store}{ui}{7.4}{guard opens the welcome screen}

  \lifelines{7.75}{ui, store, storage, db, auth}
\end{tikzpicture}

  \caption{Deleting an account. Removing the Auth user removes every row that belongs to the user, except
    the reports, which are kept without it.}
  \label{fig:seq-delete}
\end{figure}

\code{delete_account} first puts the user's email address on the ban list if the account is hidden or
suspended (\cref{sec:signup-checks}), and returns the user's chat partner, if any, to \code{open}, so
the partner can be found and can say hi again; then it deletes the Auth user. Every user foreign key
except the two in \code{reports} cascades, so the delete removes the profile, the live session and map
position, every request sent or received with its messages, every event hosted with its members and
messages, every membership, every message sent, every block made or received, the location-change log,
the event-post log and the call counts. Reports made or received stay, with the user's id set to null
(\cref{sec:blocking-reporting}). The app then
signs out and forgets what the phone saved for the account, and the Auth listener resets the store to its
initial state, closes every channel and clears the image cache.
